% GENERATED FILE. Edit canonical lessons, then run npm run generate:books.
% canonical-source-hash: fnv1a64:6c60e5314077108d
% canonical-lessons: GU-C145-bole-che, GU-C145-teo, GU-W46-o, GU-C145-bolie-chie, GU-C145-nathi-bolto, GU-C145-na-nahi, GU-C145-koi-nathi, GU-C145-not-practice

\chapter{Saying not}
\label{ch:gu-saying-not}
\hlchaptermodality{145}

\begin{chapteropening}
\textbf{By the end of this chapter:} \emph{I can complete the present for every person, say that someone does not do something with nathī, na and nahī̃, and say nobody, nothing and never.}

Turn present sentences into negatives, then carry a short exchange that needs nathī, na and nahī̃.
\end{chapteropening}

\section[\texorpdfstring{\gu{તે} … \gu{બોલે} \gu{છે} (te … bole che)}{તે … બોલે છે (te … bole che)}]{\gu{તે} \gu{ગુજરાતી} \gu{બોલે} \gu{છે} — the present, for someone else}

\label{lesson:GU-C145-bole-che}



\begin{quote}
Before the new one: say the Gujarati for \emph{to change}, the last verb of the A1 words. Then the first sentence you ever built about yourself.

\begin{itemize}
\raggedright
  \item \emph{Recall:} \emph{badalvũ}, then \emph{hũ gujarātī bolũ chũ}
\end{itemize}

That sentence has three working pieces: the stem \textbf{\gu{બોલ}-} (\emph{speak}), the ending \textbf{-\gu{ું}} that says \emph{I}, and \textbf{\gu{છું}}, the copula that agrees with \emph{I}. Change the person and two of the three pieces move.
\end{quote}

\begin{grammarlens}[title={Grammar Lens: one ending for he, she and \gu{તું}}]
\begin{quote}
\textbf{\gu{તે} \gu{ગુજરાતી} \gu{બોલે} \gu{છે}.} — \emph{te gujarātī bole che.} — He speaks Gujarati. She speaks Gujarati.
\end{quote}

The stem stays \textbf{\gu{બોલ}-}. The person ending becomes \textbf{-\gu{ે}}, and the copula becomes \textbf{\gu{છે}}, the commonest word in this book. The familiar \emph{you} takes the very same pair:

\begin{quote}
\textbf{\gu{તું} \gu{ગુજરાતી} \gu{બોલે} \gu{છે}.} — \emph{tũ gujarātī bole che.} — You speak Gujarati.
\end{quote}

\noindent
\begin{tabularx}{\linewidth}{@{}>{\raggedright\arraybackslash}X>{\raggedright\arraybackslash}X>{\raggedright\arraybackslash}X@{}}
\toprule
\textbf{who} & \textbf{verb} & \textbf{copula} \\
\midrule
\gu{હું} & \gu{બોલું} & \gu{છું} \\
\gu{તું}, \gu{તે} & \gu{બોલે} & \gu{છે} \\
\bottomrule
\end{tabularx}

\textbf{\gu{તું}} and \textbf{\gu{તે}} share one verb form, so only the pronoun in front tells them apart. \textbf{\gu{જવું}} bends its stem a little and keeps the idea: \textbf{\gu{તે} \gu{જાય} \gu{છે}}, \emph{she goes}.
\end{grammarlens}

\subsection*{Your turn}
Say these aloud:

\begin{itemize}
\raggedright
  \item \emph{hũ gujarātī bolũ chũ}, then \emph{te gujarātī bole che}
  \item \textbf{\gu{તે} \gu{વાંચે} \gu{છે}} — \emph{te vā̃ce che}, she reads
  \item \textbf{\gu{તે} \gu{અહીં} \gu{રહે} \gu{છે}} — \emph{te ahī̃ rahe che}, he lives here
  \item the last two again to a friend, with \textbf{\gu{તું}} in front
\end{itemize}

\begin{tcolorbox}[breakable,colback=teal!4,colframe=teal!35!black,title={Before you move on}]
Which two pieces change when \emph{I speak} becomes \emph{she speaks}? (The ending \textbf{-\gu{ું}} becomes \textbf{-\gu{ે}}, and \textbf{\gu{છું}} becomes \textbf{\gu{છે}}.) Which pronoun shares her verb form? (\textbf{\gu{તું}}, the familiar \emph{you}.)
\end{tcolorbox}
\section[\texorpdfstring{\gu{તેઓ} (teo)}{તેઓ (teo)}]{\gu{તેઓ} (teo) — they}

\label{lesson:GU-C145-teo}



\begin{quote}
Before the new one: say \emph{she speaks Gujarati}. Then the pronoun on its own, \textbf{\gu{તે}}, the word for anyone who is not here.
\end{quote}

\subsection*{You'll want to know: \gu{તેઓ}}
\textbf{\gu{તેઓ}} (\emph{teo}) — \textquotedblleft{}they\textquotedblright{}.

It is \textbf{\gu{તે}} with \textbf{-\gu{ઓ}} on the end. \textbf{\gu{તે}} names one person who is not here; \textbf{\gu{તેઓ}} names more than one.

\begin{quote}
\textbf{\gu{તેઓ} \gu{ગુજરાતી} \gu{બોલે} \gu{છે}.} — \emph{teo gujarātī bole che.} — They speak Gujarati.
\end{quote}

The verb does not move: \textbf{\gu{બોલે} \gu{છે}} serves \textbf{\gu{તેઓ}} exactly as it serves \textbf{\gu{તે}}. In the present, \emph{he}, \emph{she}, familiar \emph{you} and \emph{they} all meet on \textbf{-\gu{ે} \gu{છે}}.

\textbf{\gu{તેઓ}} does a second job. Said of one person, it is a respectful \emph{he} or \emph{she}: a teacher, a grandparent, a guest. So \textbf{\gu{તેઓ} \gu{અહીં} \gu{છે}} is \emph{they are here}, or \emph{he is here}, said with respect.

\begin{cousinweb}
\textbf{\gu{તેઓ}} (\emph{teo}) is \textbf{\gu{તે}} plus \textbf{-\gu{ઓ}}, the ending Gujarati nouns take for \textquotedblleft{}more than one\textquotedblright{}. \textbf{\gu{તે}} continues the old Sanskrit demonstrative stem \emph{ta-}, whose plural \emph{te} already meant \textquotedblleft{}they\textquotedblright{} (CDIAL 5612).
\end{cousinweb}

\begin{grammarlens}[title={What you've built}]
One more pronoun, and the present now has four people on one ending.
\end{grammarlens}

\subsection*{Your turn}
Say these aloud:

\begin{itemize}
\raggedright
  \item \emph{teo}
  \item \textbf{\gu{તેઓ} \gu{અહીં} \gu{રહે} \gu{છે}} — \emph{teo ahī̃ rahe che}, they live here
  \item \emph{te}, then \emph{teo}: one person, then more than one
\end{itemize}

\begin{tcolorbox}[breakable,colback=teal!4,colframe=teal!35!black,title={Before you move on}]
What does \textbf{\gu{તેઓ}} mean? (\emph{They}, or a respectful \emph{he} or \emph{she}.) Which letter at its end have you read but never written? (\textbf{\gu{ઓ}}, the next lesson.)
\end{tcolorbox}
\section[\texorpdfstring{\gu{ઓ} (o)}{ઓ (o)}]{\gu{ઓ} — the standing vowel o}

\label{lesson:GU-W46-o}



\begin{quote}
\begin{itemize}
\raggedright
  \item \emph{Recall:} say \emph{they} — \textbf{R1}, one lesson back
  \item \emph{Recall:} say the standing vowel \textbf{\gu{આ}}, and the sign \textbf{\gu{ો}} that gives a consonant its \emph{o}
\end{itemize}

\textbf{\gu{તેઓ}} ends on a vowel that stands alone. After a consonant, \emph{o} is a sign: \textbf{\gu{તો}}, \textbf{\gu{જો}}. At the start of a syllable there is no consonant to carry it, so \emph{o} needs a letter of its own.
\end{quote}

\subsection*{You'll want to know: \gu{આ} with an arc above}
\textbf{\gu{ઓ}} is built on \textbf{\gu{આ}}. Keep the whole of \textbf{\gu{આ}}, the body and its stems on the right, and set one high arc over the top.

\noindent
\begin{tabularx}{\linewidth}{@{}>{\raggedright\arraybackslash}X>{\raggedright\arraybackslash}X>{\raggedright\arraybackslash}X@{}}
\toprule
\textbf{letter} & \textbf{sound} & \textbf{what it is made of} \\
\midrule
\gu{આ} & ā & the body and its stems \\
\gu{ઓ} & o & \gu{આ}, then a high arc \\
\bottomrule
\end{tabularx}

\subsection*{Script — meet one form}
\begin{quote}
\textbf{\gu{ઓ}} — \emph{o}
\end{quote}

Keep \textbf{\gu{ઓ}} visible. Write the joined body of \textbf{\gu{આ}} first: the open curve on the left, the broad lower body, the small shoulder and arch. Lift, and bring the first stem down the right side to its foot. Lift again and bring the trailing stem down beside it. Lift a last time and sweep the high arc from left to right over the top. Four movements, three lifts.

The arc sits high and clear of the stems. Without it the letter reads \textbf{\gu{આ}}.

\subsection*{Writing — notice and trace}
\hlblockfigure{\detokenize{figures/GU-W46-o-filmstrip.pdf}}{How \gu{ઓ} is written, stroke by stroke}

Finger-trace \textbf{\gu{ઓ}} once, then air-trace it once. Notice the body of \textbf{\gu{આ}}, the two stems, and the arc that rides above them.

\subsection*{Writing — copy with the model}
Copy visible \textbf{\gu{ઓ}} once. Compare only two things: are both stems there, and does the arc clear their tops? Repair one landmark if needed.

\subsection*{Writing — recall after hiding}
Hide the model and write the standing vowel \emph{o} once. Uncover \textbf{\gu{ઓ}}, compare, and stop after one repair.

\begin{tcolorbox}[breakable,colback=teal!4,colframe=teal!35!black,title={Before you move on}]
Read \textbf{\gu{આ} \gu{ઓ} \gu{તેઓ}} aloud. What turns \textbf{\gu{આ}} into \textbf{\gu{ઓ}}? (\textbf{The high arc.})

Sources: \href{https://www.t30apps.com/gujarati-alphabet-writing/}{Gujarati Alphabet Writing Practice} and the \href{https://www.unicode.org/charts/PDF/U0A80.pdf}{Unicode Gujarati chart}.
\end{tcolorbox}
\section[\texorpdfstring{\gu{અમે} … \gu{બોલીએ} \gu{છીએ} (ame … bolīe chīe)}{અમે … બોલીએ છીએ (ame … bolīe chīe)}]{\gu{અમે} \gu{ગુજરાતી} \gu{બોલીએ} \gu{છીએ} — the present for we, and for you}

\label{lesson:GU-C145-bolie-chie}



\begin{quote}
Say \emph{they speak Gujarati}. Then the two words for \emph{we}: \textbf{\gu{અમે}}, which leaves the listener out, and \textbf{\gu{આપણે}}, which takes the listener in.

\begin{itemize}
\raggedright
  \item \emph{Recall:} \emph{teo gujarātī bole che}; \emph{ame}, \emph{āpṇe}
  \item \emph{Write it:} \textbf{\gu{ઓ}} once from memory
\end{itemize}
\end{quote}

\begin{grammarlens}[title={Grammar Lens: we}]
\begin{quote}
\textbf{\gu{અમે} \gu{ગુજરાતી} \gu{બોલીએ} \gu{છીએ}.} — \emph{ame gujarātī bolīe chīe.} — We speak Gujarati.
\end{quote}

The ending is \textbf{-\gu{ીએ}}, and the copula has a form you have not met before: \textbf{\gu{છીએ}} (\emph{chīe}). Both finish on the same \textbf{-\gu{ીએ}}, so the sentence chimes. \textbf{\gu{આપણે}} takes the same pair: \textbf{\gu{આપણે} \gu{ગુજરાતી} \gu{બોલીએ} \gu{છીએ}}, \emph{we (you and I) speak Gujarati}.
\end{grammarlens}

\subsection*{You'll want to know: respectful you}
\textbf{\gu{તમે} \gu{કેમ} \gu{છો}?} has carried \textbf{\gu{છો}} since the second day. The verb beside \textbf{\gu{તમે}} takes the matching \textbf{-\gu{ો}}:

\begin{quote}
\textbf{\gu{તમે} \gu{ગુજરાતી} \gu{બોલો} \gu{છો}.} — \emph{tame gujarātī bolo cho.} — You speak Gujarati.
\end{quote}

The present endings for one person:

\noindent
\begin{tabularx}{\linewidth}{@{}>{\raggedright\arraybackslash}X>{\raggedright\arraybackslash}X>{\raggedright\arraybackslash}X@{}}
\toprule
\textbf{who} & \textbf{verb} & \textbf{copula} \\
\midrule
\gu{હું} & \gu{બોલું} & \gu{છું} \\
\gu{તું}, \gu{તે} & \gu{બોલે} & \gu{છે} \\
\bottomrule
\end{tabularx}

And for more than one, or for respect:

\noindent
\begin{tabularx}{\linewidth}{@{}>{\raggedright\arraybackslash}X>{\raggedright\arraybackslash}X>{\raggedright\arraybackslash}X@{}}
\toprule
\textbf{who} & \textbf{verb} & \textbf{copula} \\
\midrule
\gu{અમે}, \gu{આપણે} & \gu{બોલીએ} & \gu{છીએ} \\
\gu{તમે} & \gu{બોલો} & \gu{છો} \\
\gu{તેઓ} & \gu{બોલે} & \gu{છે} \\
\bottomrule
\end{tabularx}

\subsection*{Your turn}
\begin{itemize}
\raggedright
  \item \emph{Say it:} \emph{ame gujarātī bolīe chīe}
  \item \emph{Say it:} \textbf{\gu{અમે} \gu{અહીં} \gu{રહીએ} \gu{છીએ}} — \emph{ame ahī̃ rahīe chīe}, we live here
  \item \emph{Say it:} \textbf{\gu{તમે} \gu{ક્યાં} \gu{રહો} \gu{છો}?} — \emph{tame kyā̃ raho cho?}, where do you live?
  \item \emph{Answer:} with \emph{hũ … rahũ chũ}
\end{itemize}

\begin{tcolorbox}[breakable,colback=teal!4,colframe=teal!35!black,title={Before you move on}]
Which copula ends on the same sound as its verb for \emph{we}? (\textbf{\gu{છીએ}}, after \textbf{-\gu{ીએ}}.) Which ending goes with \textbf{\gu{તમે}}? (\textbf{-\gu{ો}}, then \textbf{\gu{છો}}.)
\end{tcolorbox}
\section[\texorpdfstring{\gu{નથી} …\gu{તો} (nathī …-to)}{નથી …તો (nathī …-to)}]{\gu{હું} \gu{નથી} \gu{બોલતો} — saying \emph{not} about what people do}

\label{lesson:GU-C145-nathi-bolto}



\begin{quote}
Say \emph{we speak Gujarati} and \emph{you speak Gujarati}. Then the repair sentence from long ago: \textbf{\gu{સમજતો} \gu{નથી}}, \emph{I do not understand}.

That sentence already holds the whole negative. Here is how it is built.
\end{quote}

\begin{grammarlens}[title={Grammar Lens: \gu{નથી} takes the copula's place}]
\begin{quote}
\textbf{\gu{હું} \gu{ગુજરાતી} \gu{નથી} \gu{બોલતો}.} — \emph{hũ gujarātī nathī bolto.} — I don't speak Gujarati.
\end{quote}

Two changes, and only two:

\begin{enumerate}
\raggedright
  \item \textbf{\gu{છું}} goes, and \textbf{\gu{નથી}} comes in. \textbf{\gu{નથી}} is \emph{is not} worn into one word, and it is the same for every person: \textbf{\gu{હું} \gu{નથી}}, \textbf{\gu{તે} \gu{નથી}}, \textbf{\gu{અમે} \gu{નથી}}.
  \item The verb drops its person ending and takes \textbf{-\gu{તો}}, the shape \textbf{\gu{સમજતો}} has had all along.
\end{enumerate}

\textbf{-\gu{તો}} is the man's form. A woman says \textbf{\gu{હું} \gu{ગુજરાતી} \gu{નથી} \gu{બોલતી}}, with \textbf{-\gu{તી}}. A group, or a respectful \emph{you}, take \textbf{-\gu{તા}}: \textbf{\gu{અમે} \gu{નથી} \gu{બોલતા}}, \textbf{\gu{તમે} \gu{નથી} \gu{બોલતા}}.

\textbf{\gu{નથી}} may also stand after the verb, as it does in \textbf{\gu{સમજતો} \gu{નથી}}. Both orders are everyday Gujarati: \textbf{\gu{હું} \gu{ગુજરાતી} \gu{બોલતો} \gu{નથી}} says the same thing.
\end{grammarlens}

\begin{grammarlens}[title={Grammar Lens: when there is no other verb}]
With the copula alone, \textbf{\gu{નથી}} replaces \textbf{\gu{છે}} and nothing else moves: \textbf{\gu{તે} \gu{અહીં} \gu{નથી}}, \emph{she is not here}. That was the first job \textbf{\gu{નથી}} ever did for you.
\end{grammarlens}

\subsection*{Your turn}
Say these aloud:

\begin{itemize}
\raggedright
  \item \emph{hũ cā nathī pīto}, or \emph{pītī}, I don't drink tea
  \item \textbf{\gu{તે} \gu{ગુજરાતી} \gu{નથી} \gu{વાંચતી}} — \emph{te gujarātī nathī vā̃ctī}, she doesn't read Gujarati
  \item \textbf{\gu{અમે} \gu{અહીં} \gu{નથી} \gu{રહેતા}} — \emph{ame ahī̃ nathī rahetā}, we don't live here
\end{itemize}

\begin{tcolorbox}[breakable,colback=teal!4,colframe=teal!35!black,title={Before you move on}]
Which word replaces the copula? (\textbf{\gu{નથી}}, the same for every person.) What ending does the verb take? (\textbf{-\gu{તો}}, \textbf{-\gu{તી}} or \textbf{-\gu{તા}}: a man, a woman, or more than one.)
\end{tcolorbox}
\section[\texorpdfstring{\gu{ન} … · … \gu{નહીં} (na … · … nahī̃)}{ન … · … નહીં (na … · … nahī̃)}]{\gu{ન} and \gu{નહીં} — the two smaller \emph{nots}}

\label{lesson:GU-C145-na-nahi}



\begin{quote}
Say \emph{I don't drink tea}. Then the reply you have given since the first chapters to thanks or to an apology: \textbf{\gu{વાંધો} \gu{નહીં}}, \emph{no problem}. And the word for \emph{worry}: \textbf{\gu{ચિંતા}}.
\end{quote}

\begin{grammarlens}[title={Grammar Lens: \gu{ન} in front of a request}]
\textbf{\gu{નથી}} denies a statement. A request is not a statement, so it takes the bare \textbf{\gu{ન}} (\emph{na}) in front of the verb:

\begin{quote}
\textbf{\gu{ચિંતા} \gu{ન} \gu{કરો}.} — \emph{cintā na karo.} — Don't worry.
\end{quote}

\textbf{\gu{કરો}} is the polite \emph{do!} that \textbf{\gu{માફ} \gu{કરો}} has carried since the courtesy chapter. Put \textbf{\gu{ન}} in front of it and the request becomes a request \emph{not} to.
\end{grammarlens}

\subsection*{You'll want to know: \gu{નહીં}, for one word}
\textbf{\gu{નહીં}} (\emph{nahī̃}) cancels a word or a phrase rather than a whole sentence:

\begin{quote}
\textbf{\gu{ચા} \gu{નહીં}, \gu{પાણી}.} — \emph{cā nahī̃, pāṇī.} — Not tea: water. \textbf{\gu{હમણાં} \gu{નહીં}.} — \emph{hamṇā̃ nahī̃.} — Not now.
\end{quote}

\textbf{\gu{વાંધો} \gu{નહીં}} is built this way: \emph{objection — none}.

Three negatives, three jobs:

\begin{itemize}
\raggedright
  \item \textbf{\gu{નથી}} — a statement: \emph{it is not so};
  \item \textbf{\gu{ન}} — a request: \emph{don't};
  \item \textbf{\gu{નહીં}} — one word: \emph{not this}.
\end{itemize}

\subsection*{Your turn}
\begin{itemize}
\raggedright
  \item \emph{Say it:} \emph{cintā na karo}
  \item \emph{Say it:} \textbf{\gu{ચા} \gu{નહીં}, \gu{દૂધ}} — \emph{cā nahī̃, dūdh}, not tea: milk
  \item \emph{Say it:} \emph{hamṇā̃ nahī̃}
  \item \emph{Choose:} \textbf{\gu{નથી}}, \textbf{\gu{ન}} or \textbf{\gu{નહીં}} for \emph{she isn't here} (\textbf{\gu{નથી}}: a statement)
\end{itemize}

\begin{tcolorbox}[breakable,colback=teal!4,colframe=teal!35!black,title={Before you move on}]
Which negative goes in front of \textbf{\gu{કરો}}? (\textbf{\gu{ન}}.) Which one cancels a single word? (\textbf{\gu{નહીં}}.)
\end{tcolorbox}
\section[\texorpdfstring{\gu{કોઈ} … \gu{નથી} (koī … nathī)}{કોઈ … નથી (koī … nathī)}]{\gu{કોઈ} \gu{નથી} — nobody, nothing, never}

\label{lesson:GU-C145-koi-nathi}



\begin{quote}
Say \emph{don't worry}, then \emph{not now}. And two words from the chapters on this and that: \textbf{\gu{કોઈ}}, \emph{someone}, and \textbf{\gu{કંઈક}}, \emph{something}.
\end{quote}

\begin{grammarlens}[title={Grammar Lens: \gu{કોઈ} with \gu{નથી}}]
\begin{quote}
\textbf{\gu{અહીં} \gu{કોઈ} \gu{નથી}.} — \emph{ahī̃ koī nathī.} — There is nobody here.
\end{quote}

Gujarati needs no separate word for \emph{nobody}. It says \emph{someone is-not}: \textbf{\gu{કોઈ}}, then \textbf{\gu{નથી}}.
\end{grammarlens}

\subsection*{You'll want to know: \gu{કંઈ} and \gu{ક્યારેય}}
\textbf{\gu{કંઈક}} is \textbf{\gu{કંઈ}} with an extra \textbf{-\gu{ક}}. Before a negative, Gujarati uses the shorter \textbf{\gu{કંઈ}} (\emph{kaī̃}):

\begin{quote}
\textbf{\gu{હું} \gu{કંઈ} \gu{નથી} \gu{ખાતો}.} — \emph{hũ kaī̃ nathī khāto.} — I'm not eating anything.
\end{quote}

\textbf{\gu{ક્યારેય}} (\emph{kyārey}) is \textbf{\gu{ક્યારે}}, \emph{when}, with \textbf{-\gu{ય}}, \emph{ever}, on the end. With \textbf{\gu{નથી}} it means \emph{never}:

\begin{quote}
\textbf{\gu{તે} \gu{ક્યારેય} \gu{ચા} \gu{નથી} \gu{પીતી}.} — \emph{te kyārey cā nathī pītī.} — She never drinks tea.
\end{quote}

The shape is the same three times: an \emph{any} word, then \textbf{\gu{નથી}}.

\subsection*{Your turn}
Say these aloud:

\begin{itemize}
\raggedright
  \item \emph{ahī̃ koī nathī}
  \item \emph{hũ kaī̃ nathī khāto}, or \emph{khātī}
  \item \textbf{\gu{હું} \gu{ક્યારેય} \gu{દૂધ} \gu{નથી} \gu{પીતો}} — \emph{hũ kyārey dūdh nathī pīto}, I never drink milk
\end{itemize}

\begin{tcolorbox}[breakable,colback=teal!4,colframe=teal!35!black,title={Before you move on}]
What does \textbf{\gu{કોઈ} \gu{નથી}} say, word for word? (\emph{Someone is-not}: nobody.) What does \textbf{-\gu{ય}} add to \textbf{\gu{ક્યારે}}? (\emph{Ever}; with \textbf{\gu{નથી}}, \emph{never}.) What does \textbf{\gu{કંઈક}} become before a negative? (\textbf{\gu{કંઈ}}.)
\end{tcolorbox}
\section[Practice]{Saying \emph{not} — a practice}

\label{lesson:GU-C145-not-practice}



\begin{quote}
Say \emph{nobody is here}, \emph{I'm not eating anything}, and \emph{she never drinks tea}.
\end{quote}

\subsection*{Your turn: turn each sentence over}
\begin{itemize}
\raggedright
  \item \emph{Turn:} \textbf{\gu{તે} \gu{ગુજરાતી} \gu{બોલે} \gu{છે}} into \emph{te gujarātī nathī bolto}, or \emph{boltī}
  \item \emph{Turn:} \textbf{\gu{અમે} \gu{અહીં} \gu{રહીએ} \gu{છીએ}} into \emph{ame ahī̃ nathī rahetā}
  \item \emph{Turn:} \textbf{\gu{તમે} \gu{ચા} \gu{પીઓ} \gu{છો}} into \emph{tame cā nathī pītā}
\end{itemize}

\subsection*{The exchange: three negatives at work}
\begin{quote}
\textbf{— \gu{ચા}?} \textbf{— \gu{ચા} \gu{નહીં}, \gu{પાણી}. \gu{હું} \gu{ચા} \gu{નથી} \gu{પીતી}.} \textbf{— \gu{તેઓ} \gu{ક્યાં} \gu{છે}?} \textbf{— \gu{તેઓ} \gu{અહીં} \gu{નથી}. \gu{અહીં} \gu{કોઈ} \gu{નથી}.} \textbf{— \gu{ચિંતા} \gu{ન} \gu{કરો}.}
\end{quote}

Say these aloud:

\begin{itemize}
\raggedright
  \item the exchange aloud, taking each part in turn
  \item the second line as a man would say it: \emph{pīto}
\end{itemize}

\begin{tcolorbox}[breakable,colback=teal!4,colframe=teal!35!black,title={Before you move on}]
Which three negatives did the exchange use, and what job did each do? (\textbf{\gu{નહીં}} for one word, \textbf{\gu{નથી}} for a statement, \textbf{\gu{ન}} for a request.)
\end{tcolorbox}
